\documentclass[11pt]{article}

% ---------- Packages ----------
\usepackage[margin=1in]{geometry}
\usepackage{amsmath,amssymb,mathtools}
\usepackage{booktabs}
\usepackage{array}
\usepackage{enumitem}
\usepackage{xcolor}
\usepackage{listings}
\usepackage{algorithm}
\usepackage{algpseudocode}
\usepackage{fancyhdr}
\usepackage{titlesec}
\usepackage{caption}
\usepackage[hidelinks]{hyperref}
\usepackage{parskip}
\usepackage{longtable}

% ---------- Colors ----------
\definecolor{codebg}{RGB}{245,245,247}
\definecolor{darkblue}{RGB}{20,40,90}
\definecolor{accent}{RGB}{120,40,140}

% ---------- Listings style ----------
\lstset{
  backgroundcolor=\color{codebg},
  basicstyle=\ttfamily\small,
  breaklines=true,
  frame=single,
  rulecolor=\color{black!20},
  columns=fullflexible,
  showstringspaces=false,
  xleftmargin=4pt,
  xrightmargin=4pt,
  aboveskip=8pt,
  belowskip=8pt
}

% ---------- Section styling ----------
\titleformat{\section}{\normalfont\Large\bfseries\color{darkblue}}{\thesection}{0.6em}{}
\titleformat{\subsection}{\normalfont\large\bfseries\color{darkblue}}{\thesubsection}{0.6em}{}
\titleformat{\subsubsection}{\normalfont\bfseries}{\thesubsubsection}{0.6em}{}

\pagestyle{fancy}
\fancyhf{}
\fancyhead[L]{\small Evidence-Driven Dataset Intelligence System --- Research-Grounded Proposal (v2)}
\fancyhead[R]{\small\thepage}
\renewcommand{\headrulewidth}{0.4pt}

% New tag box for "NEW / REVISED" markers
\newcommand{\newtag}{\textcolor{accent}{\textbf{[NEW]}}\ }
\newcommand{\revtag}{\textcolor{accent}{\textbf{[REVISED]}}\ }

\title{\vspace{-1.5cm}\textbf{\Large Evidence-Driven Dataset Intelligence and Recommendation System}\\[4pt]
\large LLM-assisted discovery, content-aware suitability estimation, long-tail exploration,\\ provenance, calibrated confidence, and explainable ranking\\[10pt]
\normalsize\textit{Research-Grounded Proposal --- Version 2 (Review-Integrated Revision)}}
\author{Pratham Kashyap \\ \small Reg.\ No.\ 23BAI10455 \\ \small B.Tech.\ CSE (AI/ML), School of Computer Science and Engineering (SCAI) \\ \small VIT Bhopal University}
\date{September 2026}

\begin{document}
\maketitle
\thispagestyle{empty}

\begin{center}
\fbox{\parbox{0.92\textwidth}{\small
\textbf{Status of this document.} This is Version 2 of the Capstone Project 1 research-grounded proposal.
It incorporates a structured technical review of the Version 1 proposal and closes nine identified gaps:
(1) a formalized evidence-resolution algorithm, (2) quantitative bias and calibration metrics,
(3) an LLM zero-shot baseline, (4) a cold-start transfer-estimation fallback, (5) a human evaluation
protocol for trust and explanation quality, (6) an early compute-budget validation milestone,
(7) explicit team ownership across five to six engineers, (8) a legal/ToS compliance checklist
for automated sample acquisition, and (9) a minimum-viable CLI moved into core scope. Sections carrying
\newtag or \revtag are new or materially revised relative to Version 1; unmarked sections are retained
with light editorial changes for consistency.}}
\end{center}

\vspace{0.3cm}
\tableofcontents
\newpage

% =========================================================================
\section{Executive Summary}
% =========================================================================
The project is an evidence-driven dataset intelligence system for Machine Learning (ML), Deep Learning
(DL), and Computer Vision (CV) practitioners. A user supplies a natural-language description of a project.
The system converts the description into a structured task specification, discovers candidate datasets
across multiple public repositories, gathers evidence about those datasets, performs lightweight analysis
of accessible samples where feasible, estimates task-specific utility, mitigates popularity bias, and
produces a diversified, explainable recommendation set.

Version~1 of this proposal established that a metadata-plus-retrieval recommender is, by itself,
insufficiently differentiated: dataset recommendation is already a studied information-retrieval problem,
and a trained bi-encoder retriever (DataFinder) already demonstrates strong gains over keyword search
\cite{datafinder}. Version~1 therefore repositioned the project around \emph{evidence-backed selection}
rather than retrieval alone, introducing an evidence ledger, content-aware fingerprinting, task-utility
probing, and popularity-bias mitigation.

This revision (Version~2) closes the remaining gap between that design intent and a defensible,
gradable, and potentially publishable research artifact. The single highest-leverage addition is a
\textbf{formal evidence-resolution algorithm} (Section~\ref{sec:resolver}) that derives the system's
trust labels (\emph{Verified}, \emph{Corroborated}, \emph{Single-source claim}, \emph{Conflicting},
\emph{Unknown}) from a Dempster--Shafer belief-combination rule over source reliabilities, rather than
assigning them heuristically. This is the part of the system that is not a re-application of DataFinder,
DataPerf, Cleanlab, Croissant, or GIST --- it is new, mathematically specified, and independently
evaluable. Around this core, Version~2 adds quantitative bias and calibration metrics
(Section~\ref{sec:bias-calibration}), a cold-start fallback for datasets that cannot be sampled
(Section~\ref{sec:coldstart}), a human trust-evaluation protocol (Section~\ref{sec:human-eval}), an
LLM zero-shot baseline (Section~\ref{sec:ablation}), and practical project-management additions
(Sections~\ref{sec:m0} and~\ref{sec:raci}) needed to execute the plan with a five-to-six-person team on
Apple Silicon hardware.

\textbf{Central thesis (unchanged).} A dataset recommender should answer not only ``What datasets look
relevant?'' but ``Which datasets are most defensibly useful for this task, what evidence supports that
judgment, how uncertain is it, and what useful alternatives would otherwise be missed?''

% =========================================================================
\section{Project Foundation Retained from Version 1}
% =========================================================================
The system continues to target ML/DL/CV users, accepts natural-language project descriptions, operates
across multiple public dataset repositories (Kaggle, Hugging Face Hub, UCI, OpenML, and others), avoids
requiring complete dataset downloads in the default path, and returns a ranked, explained set of
recommendations rather than an unranked search-result list.

% =========================================================================
\section{Research Gap and Problem Definition}
% =========================================================================

\subsection{Why a metadata-only recommender is insufficient}
Dataset recommendation is already established as an information-retrieval problem. DataFinder constructs
a benchmark from research needs and evaluates BM25, nearest-neighbor, and dense bi-encoder retrieval,
reporting substantial gains from a trained bi-encoder over generic search approaches~\cite{datafinder}.
A project whose main novelty is ``LLM extraction + vector retrieval + metadata scoring'' would risk
reproducing an existing direction. The system's role is therefore shifted from pure retrieval toward
evidence-backed selection: metadata supports fast candidate generation and constraint filtering, while
additional evidence is drawn from dataset structure, documentation, provenance, and small accessible
samples before stronger suitability claims are made.

\subsection{Why repository metadata cannot be treated as ground truth}
Repository metadata is publisher-supplied. Documentation quality is uneven, and dataset visibility can
itself correlate with documentation quality. Licensing is especially sensitive: public reporting on
dataset licensing quality documents substantial rates of missing or incorrectly stated licenses across
widely used public ML datasets~\cite{licensing-audit}. The system therefore distinguishes a repository
\emph{claim} from a verified or cross-supported \emph{fact}. ``License: CC~BY~4.0'' may be reported as a
repository claim; it must not be silently translated into ``commercially safe.'' \revtag Version~2 makes
this distinction computable rather than descriptive --- see Section~\ref{sec:resolver}.

\subsection{Why popularity must not define suitability}
Downloads, stars, ratings, and citations are useful community evidence but are not equivalent to task
fitness. A popularity-driven ranking favors datasets with stronger ecosystem visibility and richer
metadata, creating a feedback loop in which well-known datasets become increasingly discoverable while
specialized or smaller datasets remain underexposed. The final ranking treats popularity as contextual
evidence, not core quality; a long-tail exploration mechanism searches for relevant candidates that are
not dominant on popularity signals, and the final set is diversified so that several near-identical
popular datasets do not crowd out useful alternatives. \revtag Version~2 requires this claim to be
measured, not merely designed for --- see Section~\ref{sec:bias-calibration}.

% =========================================================================
\section{Research Questions and Objectives}
% =========================================================================
\begin{enumerate}[leftmargin=*]
\item Build a normalized multi-source dataset corpus with stable identifiers and provenance pointers.
\item Extract structured task requirements using an LLM or a deterministic parser constrained by a schema.
\item Generate candidates through hybrid retrieval and hard-constraint filtering.
\item Assemble an evidence ledger for metadata, documentation, provenance, license claims, repository
      signals, and sample findings, resolved by a \textbf{formal combination rule} (\newtag
      Section~\ref{sec:resolver}) rather than by heuristic labeling.
\item Compute content-aware dataset fingerprints from small accessible samples when technically and
      legally feasible, with an explicit \textbf{transfer-estimate fallback} (\newtag
      Section~\ref{sec:coldstart}) when sampling is infeasible.
\item Estimate task-specific utility using lightweight probes rather than full model training.
\item Rank candidates with an evidence-aware, uncertainty-aware scoring model and diversify the final set,
      reporting the resulting \textbf{exposure disparity} (\newtag Section~\ref{sec:bias-calibration})
      relative to a non-diversified baseline.
\item Produce explanations that separate claims, measured evidence, assumptions, unknowns, and
      limitations, and validate that these explanations measurably change user trust (\newtag
      Section~\ref{sec:human-eval}).
\item Evaluate the system against strong retrieval baselines --- including a zero-shot LLM baseline
      (\newtag Section~\ref{sec:ablation}) --- robustness tests, calibration analysis, and a
      purpose-built long-tail benchmark.
\end{enumerate}

% =========================================================================
\section{Final Design Principles}
% =========================================================================
\begin{itemize}[leftmargin=*]
\item \textbf{Evidence over assertion:} a recommendation factor must be traceable to a source,
      observation, computed metric, or explicit assumption.
\item \textbf{Suitability over popularity:} popularity can inform confidence or community evidence but
      should not dominate task-specific utility.
\item \textbf{Long-tail discovery is intentional:} the system must preserve a candidate path for relevant
      datasets with weak visibility.
\item \textbf{Content when available:} actual sample evidence supplements repository descriptions instead
      of requiring complete downloads.
\item \textbf{Hard constraints are filters:} mandatory requirements (modality, minimum size, usage
      constraints) are not merely score penalties.
\item \textbf{Unknown is a valid state:} missing evidence lowers confidence rather than silently becoming
      a neutral or favorable value.
\item \textbf{Recommendations are sets, not only rankings:} the top five should be useful together and
      avoid redundant dataset families.
\item \textbf{Local-first and reproducible:} the core pipeline runs with open-source components and
      modest local compute, with optional acceleration where available.
\item \newtag \textbf{Confidence must be calibrated, not decorative:} an ``Evidence strength: High'' label
      is only meaningful if it is empirically shown to correlate with correctness. See
      Section~\ref{sec:bias-calibration}.
\end{itemize}

% =========================================================================
\section{Final System Architecture}
% =========================================================================
\begin{lstlisting}[basicstyle=\ttfamily\footnotesize]
USER PROJECT DESCRIPTION
        |
        v
[Requirement Analyst]            LLM / structured parser
        |
        v
[Task Specification]             task, domain, modality, target, labels, constraints
        |
        +-----------------------------+
        |                             |
        v                             v
[Hard Constraint Filter]      [Semantic Query Expansion]
        |                             |
        +---------------+-------------+
                        v
              [Hybrid Candidate Retrieval]      BM25 + dense embeddings
                        |
                        v
              [Candidate Pool: 30-100]
                        |
             +----------+-----------+
             |                      |
             v                      v
       [Evidence Layer]      [Content Evidence Layer]
       metadata/docs         sample access when possible
       provenance/license    fingerprints/quality/probes
             |                      |
             +----------+-----------+
                        v
      [Evidence Resolver: Dempster-Shafer Fusion]      NEW - Section 9.2
      -> Verified / Corroborated / Single-source /
         Conflicting / Unknown, with belief intervals
                        v
              [Task Utility Scoring]        (+ Cold-start transfer estimate, Section 12.1)
                        |
                        v
          [Long-Tail Exploration Pool]
                        |
                        v
          [Diversified Multi-Objective Reranking]   (+ Exposure-disparity audit, Section 13.1)
                        |
                        v
            [FINAL RECOMMENDATION SET]
     best overall | quality | efficient | hidden gem | alternative
                        |
                        v
              [Evidence-backed, Calibration-checked Report]
\end{lstlisting}

\subsection{Candidate-pool strategy}
The system retrieves a broader candidate pool than the final recommendation set. A recommended default is
50--100 candidates after retrieval and hard filtering, followed by deeper evidence analysis on a smaller
subset. This keeps the expensive content-aware stages bounded, and its feasibility is validated explicitly
in Milestone~M0 (Section~\ref{sec:m0}) before the team commits engineering effort around this number.

% =========================================================================
\section{Canonical Dataset Representation}
% =========================================================================
The dataset corpus uses a canonical internal record so that Kaggle, Hugging Face, UCI, OpenML, and other
repositories can be compared without pretending that their native schemas are identical. Croissant is used
where available as a standards-aligned metadata representation, while repository-specific fields remain
attached to their source rather than being flattened away.

% =========================================================================
\section{Evidence Ledger and Trust Model}
% =========================================================================
The evidence ledger is a core project component. Each important dataset claim is stored with its source,
observation time, evidence type, confidence, and any contradiction, preventing the explanation layer from
turning uncertain metadata into authoritative-sounding statements.

\textbf{License principle.} The system does not present automated license parsing as legal clearance. It
classifies, compares, flags conflicts, and exposes the evidence path; legal interpretation remains outside
the model's authority.

\subsection{Evidence states (Version 1 design)}
Version~1 specified five qualitative states: \emph{Verified}, \emph{Corroborated}, \emph{Single-source
claim}, \emph{Conflicting}, and \emph{Unknown}. These states were not previously derived from a stated
combination rule.

\subsection{\newtag Formal Evidence Resolution via Dempster--Shafer Fusion}
\label{sec:resolver}

This is the primary technical contribution of the project and the component most distinct from prior work
(DataFinder, DataPerf, Cleanlab, Croissant, GIST each address a different sub-problem; none formalizes
multi-source claim resolution for dataset trust). We adopt Dempster--Shafer evidence theory
\cite{shafer1976,dempster1968} because, unlike a simple weighted average, it natively represents
\emph{ignorance} as a first-class quantity distinct from disagreement --- which matches the system's
requirement that ``Unknown'' and ``Conflicting'' be different states rather than both collapsing to a
mid-range score.

\subsubsection{Frame of discernment and basic probability assignment}
For a single claim $c$ about a dataset $d$ (for example, ``license $=$ CC~BY~4.0'', or ``modality $=$
image''), define the binary frame of discernment $\Theta = \{\top, \bot\}$ (claim true / claim false).
Let $\mathcal{S}(c)$ be the set of evidence sources that address $c$ (a source that is silent on $c$
contributes nothing and is excluded). Each source $i \in \mathcal{S}(c)$ has a prior reliability
$r_i \in (0,1)$ (Table~\ref{tab:reliability}) and a stance $\sigma_i \in \{\text{assert-true},
\text{assert-false}\}$. Its basic probability assignment (BPA) over $2^\Theta$ is:
\begin{equation}
m_i(\top) = \begin{cases} r_i & \sigma_i = \text{assert-true} \\ 0 & \sigma_i = \text{assert-false}\end{cases}
\qquad
m_i(\bot) = \begin{cases} 0 & \sigma_i = \text{assert-true} \\ r_i & \sigma_i = \text{assert-false}\end{cases}
\end{equation}
\begin{equation}
m_i(\Theta) = 1 - r_i
\end{equation}
The mass placed on $\Theta$ represents the source's acknowledged unreliability: a source is never treated
as fully authoritative, so $r_i < 1$ strictly.

\begin{table}[h]
\centering
\small
\begin{tabular}{@{}llc@{}}
\toprule
\textbf{Source type} & \textbf{Example claim it can support} & \textbf{Prior $r_i$} \\
\midrule
Directly measured from sample & modality, approximate class count, format & $0.90$ \\
Peer-reviewed paper describing the dataset & task, domain, collection method & $0.85$ \\
Croissant-validated metadata record & license, schema, splits & $0.80$ \\
Official dataset card (repository-hosted) & license, size, task, modality & $0.70$ \\
Bare repository metadata field & license, tags, size & $0.60$ \\
Unverified community description (forum, README of unclear origin) & any & $0.35$ \\
\bottomrule
\end{tabular}
\caption{Default source-reliability priors. These are hand-set initial values for the prototype and are
tuned post hoc using the calibration procedure in Section~\ref{sec:bias-calibration}.}
\label{tab:reliability}
\end{table}

\subsubsection{Combination rule}
Sources are combined pairwise using Dempster's rule of combination. For two BPAs $m_1, m_2$ over $\Theta$:
\begin{equation}
m_{12}(A) = \frac{1}{1-K}\sum_{\substack{B,C \subseteq \Theta \\ B \cap C = A}} m_1(B)\,m_2(C), \qquad A \neq \emptyset
\end{equation}
\begin{equation}
K = \sum_{\substack{B,C \subseteq \Theta \\ B \cap C = \emptyset}} m_1(B)\,m_2(C)
\end{equation}
where $K$ is the \emph{conflict mass} produced by combining $m_1$ and $m_2$. Sources in $\mathcal{S}(c)$
are combined iteratively; Dempster's rule is commutative and associative, so the combination order does
not affect the final result, though numerical care (accumulating $\log(1-K)$ terms) is used to avoid
underflow when $|\mathcal{S}(c)|$ is large. After combining all sources, the system retains:
\begin{itemize}[leftmargin=*]
\item \textbf{Belief:} $\mathrm{Bel}(\top) = m(\top)$ --- the mass of evidence directly supporting the claim.
\item \textbf{Plausibility:} $\mathrm{Pl}(\top) = m(\top) + m(\Theta) = 1 - m(\bot)$ --- the upper bound
consistent with the evidence.
\item \textbf{Uncertainty interval:} $[\mathrm{Bel}(\top), \mathrm{Pl}(\top)]$.
\item \textbf{Cumulative conflict:} $K_{\text{total}}$, the total conflict mass accumulated across all
pairwise combinations, capturing how much the sources actively disagreed (as opposed to one source simply
being silent).
\end{itemize}

\subsubsection{Mapping to trust states}
The qualitative labels used throughout the system's explanations are now \emph{derived} quantities rather
than assigned labels:
\begin{align}
\text{Verified} &\iff \mathrm{Bel}(\top) \geq \tau_v \ \wedge\ K_{\text{total}} \leq \kappa_{\text{low}} \ \wedge\ |\mathcal{S}(c)| \geq 2 \\
\text{Corroborated} &\iff \tau_c \leq \mathrm{Bel}(\top) < \tau_v \ \wedge\ K_{\text{total}} \leq \kappa_{\text{low}} \ \wedge\ |\mathcal{S}(c)| \geq 2 \\
\text{Single-source claim} &\iff |\mathcal{S}(c)| = 1 \\
\text{Conflicting} &\iff K_{\text{total}} > \kappa_{\text{high}} \\
\text{Unknown} &\iff |\mathcal{S}(c)| = 0
\end{align}
with default thresholds $\tau_v = 0.85$, $\tau_c = 0.60$, $\kappa_{\text{low}} = 0.10$,
$\kappa_{\text{high}} = 0.30$ (tunable; see calibration in Section~\ref{sec:bias-calibration}).

\begin{algorithm}[h]
\caption{Evidence Resolution for a Single Claim}
\label{alg:resolver}
\begin{algorithmic}[1]
\Function{ResolveClaim}{$c$, sources $\mathcal{S}(c)$}
  \If{$|\mathcal{S}(c)| = 0$}
    \State \Return (\textsc{Unknown}, $\mathrm{Bel}=0$, $\mathrm{Pl}=1$)
  \EndIf
  \State $m \gets$ BPA of the first source in $\mathcal{S}(c)$
  \State $K_{\text{total}} \gets 0$
  \For{each remaining source $i \in \mathcal{S}(c)$}
    \State $m_i \gets$ BPA($i$) \Comment{using $r_i$ from Table~\ref{tab:reliability}}
    \State $K \gets$ conflict mass of combining $m$ with $m_i$
    \State $m \gets$ Dempster-combine($m$, $m_i$)
    \State $K_{\text{total}} \gets K_{\text{total}} + K \cdot (1 - K_{\text{total}})$
    \Comment{running conflict, bounded in $[0,1)$}
  \EndFor
  \State $\mathrm{Bel} \gets m(\top)$; \quad $\mathrm{Pl} \gets 1 - m(\bot)$
  \If{$|\mathcal{S}(c)| = 1$}
    \State \Return (\textsc{SingleSourceClaim}, $\mathrm{Bel}$, $\mathrm{Pl}$)
  \ElsIf{$K_{\text{total}} > \kappa_{\text{high}}$}
    \State \Return (\textsc{Conflicting}, $\mathrm{Bel}$, $\mathrm{Pl}$)
  \ElsIf{$\mathrm{Bel} \geq \tau_v$ \textbf{and} $K_{\text{total}} \leq \kappa_{\text{low}}$}
    \State \Return (\textsc{Verified}, $\mathrm{Bel}$, $\mathrm{Pl}$)
  \ElsIf{$\mathrm{Bel} \geq \tau_c$ \textbf{and} $K_{\text{total}} \leq \kappa_{\text{low}}$}
    \State \Return (\textsc{Corroborated}, $\mathrm{Bel}$, $\mathrm{Pl}$)
  \Else
    \State \Return (\textsc{SingleSourceClaim}, $\mathrm{Bel}$, $\mathrm{Pl}$) \Comment{low-confidence fallback label}
  \EndIf
\EndFunction
\end{algorithmic}
\end{algorithm}

A fully worked numeric example (three sources disagreeing on a license claim) is given in
Appendix~\ref{app:worked-example}. This algorithm feeds directly into the Evidence term $\gamma \cdot
\mathrm{Evidence}$ of the ranking model in Section~\ref{sec:ranking-model}: \textbf{Evidence}$(d)$ is
defined as the mean $\mathrm{Bel}(\top)$ across all claims resolved for $d$, penalized by mean
$K_{\text{total}}$, giving a single computed number rather than a qualitative impression.

% =========================================================================
\section{Content-Aware Dataset Fingerprinting}
% =========================================================================
When a repository permits lightweight sample access, the system creates a Dataset Fingerprint: not a
single vector claiming to summarize the entire dataset, but a compact evidence package describing what was
actually observed in the inspected sample. Sample size is adaptive rather than fixed --- a small initial
sample establishes feasibility, and the system expands the sample only when a signal is unstable and the
additional computation is justified.

% =========================================================================
\section{Data Quality Analysis}
% =========================================================================
Cleanlab-style model-assisted checks are included as a tool, not as the definition of correctness.
Cleanlab's confident-learning workflow requires suitable out-of-sample model probabilities and can
identify examples with suspicious labels~\cite{cleanlab}. These outputs are aggregated into evidence
features such as estimated label-issue rate or the share of high-risk examples.

\begin{enumerate}[leftmargin=*]
\item Obtain a task-appropriate sample and preserve the dataset's original labels.
\item Create out-of-sample prediction probabilities using cross-validation or a held-out procedure.
\item Run label-issue detection and record suspicious-example counts rather than silently deleting data.
\item Compute duplicate and near-duplicate statistics using hashes and embedding similarity.
\item Compute modality-specific structural diagnostics.
\item Store all measurements in the evidence ledger with sample size and procedure, feeding the
      \textsc{DirectlyMeasured} source type in Table~\ref{tab:reliability}.
\item Use the resulting features as quality evidence with uncertainty bounds, not as ground truth.
\end{enumerate}

% =========================================================================
\section{Task-Specific Utility Estimation}
% =========================================================================
The key extension beyond conventional dataset recommendation is a lightweight estimate of whether a
dataset appears useful for the stated ML task. DataPerf demonstrates the value of evaluating data
selection through downstream task performance~\cite{dataperf}; the system adapts that principle under
local-compute constraints. A probe score is an estimate under a chosen probe model, sample, and evaluation
split --- it is not a guarantee that the full dataset will produce the same downstream performance.

\subsection{\newtag Cold-Start Fallback via Transfer Estimation}
\label{sec:coldstart}
Section~17.5 (Sample Acquisition Procedure) already triages several sampling strategies, but Version~1 did
not specify a behavior when \emph{none} succeed --- a common outcome for paywalled, gated, or
terms-restricted repositories. Rather than silently dropping such datasets from utility scoring (which
would bias the system back toward easily-scraped, and therefore already-popular, datasets), Version~2
defines a transfer estimate.

For an unsampled candidate $d$, let $\mathcal{N}_k(d)$ be its $k$ nearest already-probed datasets in
metadata/description embedding space (using the same encoder as retrieval). Define an RBF similarity
weight $w(d,d') = \exp\!\big(-\lVert e(d)-e(d')\rVert^2 / (2\sigma^2)\big)$ and estimate:
\begin{equation}
\hat{u}(d) = \frac{\sum_{d' \in \mathcal{N}_k(d)} w(d,d') \cdot u(d')}{\sum_{d' \in \mathcal{N}_k(d)} w(d,d')}
\end{equation}
with an inflated uncertainty band that widens as the nearest probed neighbor becomes less similar:
\begin{equation}
\sigma_{\hat{u}}(d) = \sigma_{\text{base}} + \lambda \left(1 - \max_{d' \in \mathcal{N}_k(d)} w(d,d')\right)
\end{equation}
Transfer estimates are explicitly tagged \textsc{TransferredEstimate---NotDirectlyObserved} in the
evidence ledger and the explanation card, are excluded from ever producing a \textsc{Verified} content
label, and are down-weighted (not discarded) in the final ranking. This keeps restricted-but-relevant
datasets in contention without overstating what was actually measured.

% =========================================================================
\section{Popularity Bias and Long-Tail Discovery}
% =========================================================================
Popularity is separated from suitability. The retrieval stage reserves candidate slots for datasets that
meet task requirements but are not dominant on repository popularity signals; this exploration pool is
then evaluated by the same evidence and utility criteria as mainstream candidates. The objective is not to
reward obscurity --- a lesser-known dataset reaches the final recommendation only when its evidence
supports utility, diversity, efficiency, specialization, or another user-relevant advantage.
\begin{lstlisting}[basicstyle=\ttfamily\footnotesize]
candidate_score = task_fit + evidence_quality + content_utility + feasibility
long_tail_pool  = relevant_candidates
                  filtered to non-dominant popularity percentiles
                  while preserving hard constraints
final_set = diversified_selection(main_pool + long_tail_pool)
\end{lstlisting}

\subsection{\newtag Quantifying Bias: Exposure Disparity and Reliability Calibration}
\label{sec:bias-calibration}
Version~1 asserted that diversification would reduce popularity bias and targeted ``diversity metric
improves by more than 20\%.'' A diversity increase alone does not prove that popularity bias was reduced
--- it proves that the top-$k$ set spread out. Version~2 adds two measurements that are actually tied to
the stated research claims.

\subsubsection{Exposure Disparity Index (EDI)}
Using the benchmark of Section~\ref{sec:hidden-gem}, bucket the full retrieved candidate pool for each
query into popularity deciles $j \in \{1,\dots,10\}$ (by downloads/stars, log-scaled). Let
$R_j$ be the fraction of \emph{human-judged-relevant} candidates falling in decile $j$, and let $E_j$ be
the fraction of \emph{final recommendation slots} (aggregated over the benchmark) drawn from decile $j$.
Define:
\begin{equation}
\mathrm{EDI} = \frac{1}{2}\sum_{j=1}^{10} \left| E_j - R_j \right|
\end{equation}
This is the total variation distance between where recommendations actually land and where relevance says
they should land. $\mathrm{EDI} = 0$ means exposure exactly tracks relevance, regardless of popularity;
$\mathrm{EDI} \to 1$ means exposure is maximally distorted by popularity. The system reports
$\mathrm{EDI}$ for (a) bi-encoder-only ranking, and (b) the full system with long-tail exploration and
diversified reranking, and the hypothesis is $\mathrm{EDI}_{\text{full}} < \mathrm{EDI}_{\text{baseline}}$.
A secondary sanity-check statistic, the Gini coefficient of the popularity-decile distribution of
recommended items, is reported alongside EDI but is not the primary claim, since a lower Gini coefficient
alone does not distinguish ``spread evenly across relevant items'' from ``spread evenly regardless of
relevance.''

\subsubsection{Calibration of confidence labels}
\label{sec:calibration}
The design principle ``confidence must be calibrated, not decorative'' is tested using standard
calibration diagnostics against the benchmark's human relevance/quality judgments. Predicted confidence
$\mathrm{Bel}(\top)$ values (Section~\ref{sec:resolver}) are grouped into $B$ equal-width bins; within each
bin $b$, $\mathrm{acc}(b)$ is the observed fraction of claims confirmed correct by the benchmark's ground
truth, and $\mathrm{conf}(b)$ is the mean predicted belief in that bin. Expected Calibration Error:
\begin{equation}
\mathrm{ECE} = \sum_{b=1}^{B} \frac{n_b}{N} \left| \mathrm{acc}(b) - \mathrm{conf}(b) \right|
\end{equation}
and, for continuous utility scores $p_i \in [0,1]$ against binary outcomes $o_i$ (dataset judged suitable
or not by a probe/human), the Brier score:
\begin{equation}
\mathrm{BS} = \frac{1}{N}\sum_{i=1}^{N} (p_i - o_i)^2
\end{equation}
Both are reported as reliability diagrams (predicted confidence vs.\ observed accuracy) in the final
report, following standard practice for calibration evaluation of predictive models~\cite{guo-calibration}.
Low ECE and low Brier score are required before ``Evidence strength: High'' is allowed to appear in any
user-facing explanation card; if calibration is poor, the source-reliability priors in
Table~\ref{tab:reliability} are the first tuning target.

% =========================================================================
\section{Diversified Recommendation with Utility and Diversity}
% =========================================================================
GIST formulates a max-min diversification problem that balances submodular utility with diversity and
provides a $1/2$-approximation guarantee for the stated optimization problem~\cite{gist}. The project
adopts the principle rather than requiring full-scale implementation over millions of items. After
candidate analysis, the system operates on a bounded pool and uses a GIST-inspired greedy procedure or a
computationally lighter farthest-first/submodular alternative. Dataset-family similarity is considered in
diversification: mirrors, versions, and closely derived datasets should not occupy multiple final slots
when they offer effectively the same evidence and coverage.

% =========================================================================
\section{Final Multi-Objective Recommendation Model}
\label{sec:ranking-model}
% =========================================================================
The final recommendation model is multi-objective rather than a single metadata-weighted average:
\begin{equation}
R(D \mid Q) = \alpha \cdot \mathrm{Fit} + \beta \cdot \mathrm{Utility} + \gamma \cdot \mathrm{Evidence}
  + \delta \cdot \mathrm{Coverage} + \varepsilon \cdot \mathrm{Novelty} - \lambda \cdot \mathrm{Risk}
\end{equation}
where \textbf{Fit} captures explicit requirements, \textbf{Utility} captures content-aware task evidence
(with transfer estimates from Section~\ref{sec:coldstart} where direct probing is infeasible),
\textbf{Evidence} is now the computed quantity $\mathrm{Bel}(\top)$ averaged across claims and penalized
by $K_{\text{total}}$ (Section~\ref{sec:resolver}), \textbf{Coverage} captures diversity and
representativeness, \textbf{Novelty} rewards useful long-tail discovery, and \textbf{Risk} captures
uncertainty, contradiction, missing governance information, or operational problems.

The coefficients $(\alpha,\beta,\gamma,\delta,\varepsilon,\lambda)$ are not assumed universally correct.
Initial values are hand-designed for the prototype and subsequently tuned on the benchmark
(Section~\ref{sec:hidden-gem}) using pairwise human relevance judgments and a simple learning-to-rank
model (logistic regression over the six components, or gradient-boosted pairwise ranking if enough
judgment pairs are collected). Popularity is intentionally excluded from the core suitability term; it may
enter only as contextual evidence, long-tail stratification, or a controlled bias-audit variable (as used
directly in Section~\ref{sec:bias-calibration}).

% =========================================================================
\section{Explanation and Recommendation Output}
% =========================================================================
Each recommended dataset produces an evidence-backed recommendation card rather than a single score:
\begin{lstlisting}[basicstyle=\ttfamily\footnotesize]
Dataset: <name>
Overall recommendation: 91/100 (illustrative structure)
Best for: <task/domain/modality>
Evidence strength: High / Medium / Low      (derived from Bel(true), Section 9.2 - calibrated, Section 13.1)

Why it fits
- Task: measured / corroborated / source claim
- Domain: ...
- Modality: ...
- Labels: ...
- Observed sample quality: ...

Why it is not perfect
- Unknown or conflicting evidence: ...
- Operational limitations: ...
- [If applicable] Utility is a transferred estimate, not directly observed (Section 12.1)

Why it appears here
- More diverse than the top popular alternatives
- Smaller but cleaner / better balanced

Sources and timestamps
- Repository metadata
- Dataset documentation / paper
- Sample analysis
\end{lstlisting}
This card is the artifact evaluated directly in the human trust study (Section~\ref{sec:human-eval}).

% =========================================================================
\section{Local Implementation Strategy}
% =========================================================================
The core implementation targets ordinary personal computers, specifically the team's available hardware
(MacBook Air M1, 8\,GB RAM). The system is CPU-first, memory-conscious, cache-heavy, and modular. Optional
acceleration can be enabled when a supported local device provides it. The project does not depend on
proprietary APIs for its central pipeline.

\subsection{Environment setup procedure}
\begin{lstlisting}[language=bash]
python3 -m venv .venv
source .venv/bin/activate
python -m pip install --upgrade pip

pip install pandas numpy scipy scikit-learn
pip install sentence-transformers
pip install faiss-cpu
pip install cleanlab
pip install mlcroissant
# Optional retrieval / diversification components
pip install pyserini
pip install submodlib
\end{lstlisting}
Package versions are pinned after the first successful environment build. The repository contains a
requirements/lock file and a small environment-validation script.

\subsection{Repository ingestion procedure}
\begin{enumerate}[leftmargin=*]
\item Create one source adapter per repository; never let repository-specific JSON fields leak directly
      into downstream ranking logic.
\item Assign a stable internal dataset identifier and retain the native source identifier.
\item Normalize task, modality, format, class count, sample count, license claim, documentation links,
      timestamps, and popularity signals into the canonical schema.
\item Retain raw source snapshots or request/response files for reproducibility where licensing and
      storage terms permit (see Section~\ref{sec:legal} below).
\item Attempt Croissant parsing when a valid record is available; retain provenance links and validation
      results.
\item Generate the retrieval document from a controlled concatenation of name, description, tags, task
      semantics, and structured textual fields.
\item Precompute dense embeddings and store them in the local vector index.
\item Build a BM25 index over the same normalized textual representation.
\item Cache every source fetch so experiments do not repeatedly hit external services.
\end{enumerate}

\subsection{\newtag Legal and ToS Compliance Checklist for Automated Sample Acquisition}
\label{sec:legal}
Section~17.5's sample-acquisition ordering (streaming/range requests, sample endpoints, previews, manual
acquisition) is a technical strategy; Version~1 did not separately confirm that each strategy is
permitted by the relevant platform. Before any automated fetch of a repository, the ingestion adapter for
that source must record:
\begin{enumerate}[leftmargin=*]
\item Whether the platform's API terms of service permit automated/programmatic access for the intended
      use (research/educational, non-redistributive).
\item The platform's documented rate limits, and a corresponding client-side rate limiter (never inferred
      empirically by trial and error against the live service).
\item Whether \texttt{robots.txt} or an equivalent access policy restricts the specific endpoints used.
\item Whether the dataset's own license permits the sampling/inspection the system performs, independent
      of the platform's API terms (a dataset can be API-accessible but license-restricted).
\end{enumerate}
This checklist is stored per source adapter (not per dataset) and is reviewed once per repository
integration, not once per request. The caching-first policy already specified in Version~1 additionally
serves this compliance purpose: every source fetch is cached so experiments do not repeatedly re-request
the same resource.

\subsection{Requirement extraction procedure}
\begin{lstlisting}[basicstyle=\ttfamily\footnotesize]
{
  "task": "image classification",
  "domain": "plant disease",
  "modality": "image",
  "learning_type": "supervised",
  "label_type": "multiclass",
  "min_samples": 5000,
  "required_classes": null,
  "license_requirement": "academic_use",
  "hard_constraints": [],
  "soft_preferences": []
}
\end{lstlisting}
The schema rejects unknown fields or places them into an explicit ``unresolved constraints'' field. This
makes requirement-extraction errors measurable rather than invisible.

\subsection{Hybrid retrieval procedure}
\begin{enumerate}[leftmargin=*]
\item Run BM25 for exact terms such as disease names, modalities, target variables, file formats, or
      domain-specific vocabulary.
\item Run dense retrieval on the full natural-language and structured query representation.
\item Fuse the candidate lists using reciprocal-rank fusion or calibrated score normalization.
\item Apply hard constraints after candidate generation so recall is not unnecessarily lost during the
      first lexical/dense pass.
\item Deduplicate mirrors and obvious dataset-family duplicates before deeper analysis.
\item Retain 30--100 candidates for the evidence stage and log the retrieval source that contributed each
      candidate.
\end{enumerate}

\subsection{Sample acquisition procedure}
The system avoids indiscriminate full downloads. Preferred order: direct streaming or repository viewer
APIs; small file/range requests where supported; public sample endpoints; repository-provided previews;
finally, limited manual acquisition for benchmark datasets. Every sample-analysis record states the number
of observations inspected and the acquisition method, and is checked against
Section~\ref{sec:legal} before use.

\subsection{Probe computation procedure}
\begin{enumerate}[leftmargin=*]
\item Feasibility check: can the candidate be sampled within a bounded size and time budget
      (validated empirically in Milestone~M0, Section~\ref{sec:m0})?
\item If infeasible, invoke the transfer-estimate fallback (Section~\ref{sec:coldstart}) instead of
      dropping the candidate.
\item Create a task-appropriate representation: image embeddings, text embeddings, or normalized tabular
      features.
\item Compute structural and distributional diagnostics.
\item If labels are available, train a small probe model with stratified validation and obtain
      out-of-sample predictions.
\item Run Cleanlab-style label-issue estimation where assumptions are satisfied.
\item Compute class coverage and representation diversity.
\item Persist all intermediate metrics so ranking can be rerun without recomputing samples.
\end{enumerate}

\subsection{Provenance and license workflow}
\begin{lstlisting}[basicstyle=\ttfamily\footnotesize]
Repository claim ----+
Dataset card --------+--> Evidence Resolver (Alg. 1) --> state: verified / corroborated /
Paper / source ------+                                    single-source / conflicting / unknown
Croissant ------------+                                    with Bel(true), Pl(true), K_total
\end{lstlisting}
Automated license parsing remains a classification aid, not legal advice. Croissant~1.1 adds
machine-actionable provenance concepts, vocabulary links, and usage-policy mechanisms~\cite{croissant11};
the implementation adopts these structures where present and retains raw provenance references rather than
flattening everything into a single ``reliability score.''

% =========================================================================
\section{End-to-End Recommendation Algorithm}
% =========================================================================
\begin{lstlisting}[language=Python]
def recommend(query, k=5):
    spec = extract_requirements(query)

    lexical = bm25_retrieve(spec.text_query, top_n=100)
    dense = dense_retrieve(spec.embedding_query, top_n=100)
    candidates = fuse_and_deduplicate(lexical, dense)
    candidates = hard_filter(candidates, spec.hard_constraints)

    raw_evidence = collect_evidence(candidates)               # per-claim sources
    evidence = {d: resolve_claims(raw_evidence[d]) for d in candidates}   # NEW: Algorithm 1

    feasible = select_sample_accessible(candidates, budget=probe_budget)
    content = analyze_samples(feasible, spec)
    unsampled = [d for d in candidates if d not in feasible]
    transferred = transfer_estimate(unsampled, content)       # NEW: Section 12.1

    for d in candidates:
        d.fit = requirement_fit(d, spec)
        d.utility = content.get(d.id, transferred.get(d.id))
        d.evidence = mean_belief(evidence[d.id])               # from Algorithm 1
        d.risk = uncertainty_and_conflict(evidence[d.id])
        d.novelty = long_tail_value(d)
        d.coverage = coverage_value(d, spec)

    pool = multi_objective_rank(candidates)          # R(D|Q), Section 15
    main, long_tail = split_exploration_pool(pool)
    final = diversified_select(main + long_tail, k=k)
    audit_exposure_disparity(pool, final)             # NEW: Section 13.1, logged not blocking
    return explain_with_evidence(final, evidence, content)
\end{lstlisting}

% =========================================================================
\section{Evaluation Framework}
% =========================================================================
Evaluation is treated as a first-class research artifact. The system is not considered successful merely
because recommendations appear plausible; it must demonstrate where each additional component changes
retrieval, ranking, discovery, or trustworthiness.

\subsection{Hidden-Gem benchmark}
\label{sec:hidden-gem}
A dedicated benchmark contains queries for which expert-approved datasets include at least one
lesser-known candidate. The benchmark records relevance judgments and popularity strata before the final
system is run. A hidden gem is not defined as ``low downloads''; it is a dataset with useful relevance or
task evidence whose visibility is lower than dominant alternatives. This benchmark is also the ground
truth used for the EDI and calibration measurements in Section~\ref{sec:bias-calibration}.

\subsection{Robustness benchmark}
Synthetic corruption tests remove or contradict selected metadata fields. For example, class counts are
removed, licenses are marked unknown, and descriptions are weakened. The system should not collapse
because one metadata source is incomplete; it should lower confidence, rely on other evidence where
available, and expose uncertainty. With the formal resolver (Section~\ref{sec:resolver}), this becomes a
directly testable property: removing a corroborating source should shift a claim from
\textsc{Verified} to \textsc{Corroborated} or \textsc{Single-source claim} in a predictable direction,
which can be asserted as a unit test on Algorithm~\ref{alg:resolver} rather than only observed
qualitatively.

\subsection{Avoiding fabricated result claims}
No performance values are assumed in advance. Any target (e.g., a percentage gain) is a research
hypothesis, not a result. Final tables are populated only after experiments are run. The project reports
confidence intervals or repeated-run variability where appropriate.

\subsection{\newtag Human Evaluation of Trust and Explanation Quality}
\label{sec:human-eval}
The system's central thesis is a claim about \emph{defensibility and trust}, which automated IR and
diversity metrics cannot directly validate. A small human study is added as a required, not optional,
evaluation component.

\textbf{Design.} Recruit 10--15 participants (peers/labmates with ML background). For a fixed set of
8--10 query--dataset pairs drawn from the benchmark, each participant sees, in randomized order: (a) a
plain metadata card (name, description, raw license field, download count), and (b) the evidence-backed
recommendation card from Section~\ref{sec:resolver}/12 for the same dataset. For each, participants rate
on a 5-point Likert scale: (i) ``I would trust this recommendation enough to start using the dataset,''
and (ii) ``I understand why this dataset was or was not a good fit.'' Participants are also asked a
forced-choice question: ``which card would you act on faster?''

\textbf{Analysis.} Report mean Likert scores per card type with a paired comparison (e.g., Wilcoxon
signed-rank test given the small, non-normal sample), and the forced-choice win rate. This directly tests
the ``evidence over assertion'' design principle rather than asserting it.

% =========================================================================
\section{Baselines and Ablation Plan}
\label{sec:ablation}
% =========================================================================
\begin{enumerate}[leftmargin=*]
\item \newtag \textbf{LLM zero-shot recommendation} --- ask a general-purpose LLM directly for dataset
      recommendations with no retrieval, evidence, or probing. This is the practical alternative most
      users would otherwise reach for, and the project must show measurable gains over it (particularly
      on hidden-gem recall, license correctness, and hallucination rate --- does the LLM recommend
      datasets that do not exist or misstate their license?).
\item BM25-only retrieval.
\item Dense embedding-only retrieval.
\item DataFinder-style trained bi-encoder retrieval.
\item Hybrid retrieval without evidence scoring.
\item Hybrid + metadata suitability scoring.
\item Hybrid + evidence ledger + Dempster--Shafer resolution, no uncertainty-aware ranking.
\item Hybrid + content-aware utility estimation (including transfer estimates).
\item Full system without diversification.
\item Full system with long-tail exploration and diversified reranking.
\end{enumerate}
This sequence isolates whether each design layer contributes independently and prevents the final system
from being presented as a monolithic black box.

% =========================================================================
\section{Scope and Development Roadmap}
% =========================================================================

\subsection{Scope control}
\begin{itemize}[leftmargin=*]
\item The core MVP does not require training a large language model.
\item The core MVP does not require complete downloads of every candidate dataset.
\item Bi-encoder fine-tuning is optional until baseline retrieval is established.
\item Full GIST reproduction at web-scale is not required; the project uses a bounded implementation or an
      efficient approximation.
\item A single modality is used for initial content-probing validation, then extended to a second
      modality for research completeness.
\item \revtag A minimal command-line interface that prints the evidence-backed recommendation card is
      moved into core MVP scope (previously an optional extension). With five to six engineers available,
      its cost is small relative to its value for demonstration, grading, and the human evaluation study
      in Section~\ref{sec:human-eval}, which needs a real artifact for participants to react to.
\item The final system prefers reproducibility and measurable gains over adding many loosely integrated
      features.
\end{itemize}

\subsection{\newtag Milestone M0: Compute-Budget Validation Spike}
\label{sec:m0}
Before any subsystem is built at scale, the team runs one full evidence-collection-plus-probe cycle on
approximately 10 real datasets, on the target hardware (MacBook Air M1, 8\,GB RAM), and records:
\begin{itemize}[leftmargin=*]
\item wall-clock time per candidate for retrieval, evidence resolution, sampling, and probing;
\item peak memory usage per stage;
\item failure rate of sample acquisition per repository.
\end{itemize}
These numbers directly determine whether the 50--100 candidate-pool default (Section~5.1) and the
30--100 candidate range in the retrieval procedure are realistic, \emph{before} the team commits a
milestone schedule around them. If per-candidate probe cost exceeds a set budget (e.g., 15 seconds),
the candidate pool size or probe depth is revised at this point rather than discovered during M4--M5.
M0 is scheduled for the first 1--2 days of the project, in parallel with environment setup (M1).

\subsection{\newtag Team Ownership (Five to Six Engineers)}
\label{sec:raci}
\begin{table}[h]
\centering
\small
\begin{tabular}{@{}p{3.4cm}p{6.8cm}p{2.6cm}@{}}
\toprule
\textbf{Subsystem} & \textbf{Scope} & \textbf{Consumes / Produces} \\
\midrule
Retrieval \& Ingestion & Source adapters, canonical schema, BM25/dense indices, M0 spike &
  Produces: candidate pool \\
Evidence \& Probing & Evidence Resolver (Alg.~1), Cleanlab integration, fingerprinting,
  cold-start transfer estimate & Consumes: candidate pool; Produces: resolved evidence + utility \\
Ranking \& Diversification & Multi-objective model $R(D\mid Q)$, GIST-inspired reranking,
  long-tail pool, EDI audit & Consumes: evidence + utility; Produces: final set \\
Evaluation \& Calibration & Hidden-gem/robustness benchmarks, ablations, ECE/Brier, human study
  (owns Section~\ref{sec:human-eval}) & Consumes: full pipeline outputs \\
Reporting \& CLI/UI & Explanation card rendering, minimal CLI, documentation, final report &
  Consumes: final set + evidence \\
\bottomrule
\end{tabular}
\caption{Suggested subsystem ownership for a five-to-six-person team. Interfaces between rows are explicit
data contracts (e.g., the schema of a ``resolved evidence'' record) agreed before implementation begins,
so subsystems can be built and tested independently.}
\end{table}

% =========================================================================
\section{Risks, Limitations, and Safeguards}
% =========================================================================
\begin{itemize}[leftmargin=*]
\item \textbf{Licensing risk.} Automated license parsing can be wrong or incomplete; mitigated by treating
      license fields as claims subject to Evidence Resolution (Section~\ref{sec:resolver}), never as
      legal clearance.
\item \textbf{Privacy risk.} Sample inspection could surface sensitive content in poorly documented
      datasets; mitigated by restricting to publicly listed datasets and sanitizing any logged samples.
\item \textbf{Popularity feedback loop.} Left unchecked, ranking would reinforce already-popular datasets;
      mitigated by the long-tail pool and measured directly via EDI (Section~\ref{sec:bias-calibration}).
\item \newtag \textbf{Compliance risk in automated sampling.} Programmatic access to third-party platforms
      can violate API terms if unmanaged; mitigated by the per-source compliance checklist
      (Section~\ref{sec:legal}) and caching-first fetch policy.
\item \textbf{Compute risk.} Probe-based utility estimation is the most compute-intensive stage on
      constrained hardware; mitigated by the M0 validation spike (Section~\ref{sec:m0}) and the transfer
      fallback (Section~\ref{sec:coldstart}), which avoids forcing every candidate through a full probe.
\item \textbf{Overclaiming risk.} A system that reports high-sounding confidence without empirical support
      is worse than one that reports honest uncertainty; mitigated by the calibration requirement
      (Section~\ref{sec:calibration}), which gates any ``High'' confidence label on measured ECE/Brier
      performance.
\end{itemize}

% =========================================================================
\section{Expected Project Contribution}
% =========================================================================
The project's contribution is not a new language model or a generic RAG interface. It is an integrated
methodology for moving from dataset retrieval toward evidence-driven dataset selection, with a specific,
independently citable technical core: a \textbf{Dempster--Shafer evidence-resolution algorithm for
multi-source dataset trust claims} (Section~\ref{sec:resolver}), evaluated for both bias reduction
(Section~\ref{sec:bias-calibration}) and calibration, and validated against a human trust study
(Section~\ref{sec:human-eval}) rather than automated metrics alone. This combination --- a formal fusion
rule, a measured bias-reduction claim, and a human-validated trust claim --- is what distinguishes the
project from a re-application of DataFinder, Cleanlab, Croissant, and GIST, and is the part most suitable
for framing as a short paper or workshop submission beyond the capstone deliverable itself.

% =========================================================================
\section{Expected User-Facing Result}
% =========================================================================
For a query such as ``I want to build a CNN model for classifying plant diseases from leaf images; I need
a multi-class dataset,'' the system returns a recommendation set with different roles (best overall,
highest quality, most efficient, hidden gem, notable alternative). Every recommendation includes source
links, calibrated evidence strength, main reasons, limitations, and the specific factors that caused it to
outrank or differ from nearby alternatives.

% =========================================================================
\section{References and Research Sources}
% =========================================================================
\begin{thebibliography}{99}

\bibitem{datafinder} Lin, V.\ et al. ``DataFinder: Scientific Dataset Recommendation from Natural
Language.'' \textit{WSDM 2023}. \url{https://arxiv.org/abs/2305.16636}. Official implementation:
\url{https://github.com/viswavi/datafinder}.

\bibitem{dataperf} Mazumder, M.\ et al. ``DataPerf: Benchmarks for Data-Centric AI Development.''
\textit{arXiv:2207.10062}. Official repository: \url{https://github.com/mlcommons/dataperf}.

\bibitem{cleanlab} Northcutt, C.\ et al. ``Confident Learning: Estimating Uncertainty in Dataset
Labels.'' \textit{NeurIPS 2021}. Implementation and documentation: \url{https://github.com/cleanlab/cleanlab},
\url{https://docs.cleanlab.ai/}.

\bibitem{croissant11} MLCommons. Croissant Dataset Metadata Format Specification, including the
Croissant~1.1 provenance and usage-policy extensions. \url{https://docs.mlcommons.org/croissant/}.

\bibitem{gist} Zadimoghaddam, M.\ et al. ``GIST: Greedy Independent Set Thresholding for Max-Min
Diversification with Submodular Utility.'' \textit{arXiv:2405.18754}, NeurIPS 2025.
\url{https://arxiv.org/abs/2405.18754}.

\bibitem{gist-blog} Google Research. ``Introducing GIST: The Next Stage in Smart Sampling.''
\url{https://research.google/blog/introducing-gist-the-next-stage-in-smart-sampling/}.

\bibitem{licensing-audit} Public reporting and the Data Provenance Initiative on dataset licensing and
provenance quality in widely used public ML datasets, used to justify verification and uncertainty
rather than treating repository license fields as ground truth.

\bibitem{dataset-search} Google Dataset Search research and dataset discovery literature, used as
contextual background for large-scale dataset discovery and heterogeneous source indexing.

\bibitem{hf-hub} Hugging Face Hub dataset documentation, Dataset Cards, and dataset viewer
documentation, used as practical evidence-source patterns for dataset context and sample inspection.

\bibitem{shafer1976} Shafer, G. \textit{A Mathematical Theory of Evidence.} Princeton University Press,
1976.

\bibitem{dempster1968} Dempster, A.\ P. ``A Generalization of Bayesian Inference.'' \textit{Journal of the
Royal Statistical Society, Series B}, 1968.

\bibitem{guo-calibration} Guo, C.\ et al. ``On Calibration of Modern Neural Networks.'' \textit{ICML
2017}. Used as the methodological basis for the reliability-diagram and Expected Calibration Error
evaluation in Section~\ref{sec:calibration}.

\end{thebibliography}

\newpage
\appendix
\section*{Appendices}
\addcontentsline{toc}{section}{Appendices}

% =========================================================================
\section{Practical Repository Structure}
% =========================================================================
\begin{lstlisting}[basicstyle=\ttfamily\footnotesize]
capstone-dataset-intelligence/
|-- data/
|   |-- raw/                 # source snapshots where permitted
|   |-- normalized/          # canonical dataset records
|   |-- evidence/            # evidence ledger records (incl. Bel/Pl/K_total per claim)
|   |-- fingerprints/        # cached sample-level features
|   `-- benchmarks/          # queries + relevance judgments (hidden-gem, robustness)
|-- indexes/
|   |-- bm25/
|   `-- dense/
|-- src/
|   |-- ingestion/           # per-source adapters + ToS/compliance metadata (Sec. 17.2)
|   |-- normalization/
|   |-- retrieval/
|   |-- evidence/            # Dempster-Shafer resolver (Algorithm 1)
|   |-- probing/             # incl. cold-start transfer estimator (Sec. 12.1)
|   |-- ranking/
|   |-- diversification/
|   `-- reporting/           # explanation cards + CLI (moved into MVP scope)
|-- experiments/
|   |-- baselines/           # incl. LLM zero-shot baseline (Sec. 20)
|   |-- ablations/
|   |-- robustness/
|   |-- calibration/         # NEW: ECE / Brier / reliability diagrams
|   |-- bias_audit/          # NEW: EDI / Gini
|   `-- long_tail/
|-- human_study/             # NEW: protocol, consent notes, Likert data (Sec. 19.4)
|-- configs/
|-- scripts/
|-- tests/                   # incl. resolver unit tests (robustness corruption cases)
|-- requirements.txt
`-- README.md
\end{lstlisting}

% =========================================================================
\section{Minimum Viable System vs.\ Research System}
% =========================================================================
\begin{table}[h]
\centering
\small
\begin{tabular}{@{}p{5.2cm}p{5.2cm}@{}}
\toprule
\textbf{Minimum Viable System (M1--M6)} & \textbf{Full Research System (M7--M8+)} \\
\midrule
Single modality, hybrid retrieval, metadata evidence, basic Dempster--Shafer resolution with default
priors, greedy diversification, CLI output. &
Second modality, tuned reliability priors via calibration, full ablation suite including LLM baseline,
EDI/Gini bias audit, human trust study, GIST-inspired reranking, tuned $R(D\mid Q)$ coefficients. \\
\bottomrule
\end{tabular}
\end{table}

% =========================================================================
\section{Implementation Stop Conditions}
% =========================================================================
\begin{itemize}[leftmargin=*]
\item If dense retrieval does not beat the lexical baseline, stop tuning downstream ranking and revisit
      corpus normalization/query formulation.
\item If sample acquisition is unreliable across repositories, keep metadata/provenance evidence as the
      fallback and clearly label content evidence as unavailable (or transferred, Section~\ref{sec:coldstart}).
\item If probe models consume excessive local compute, reduce the candidate pool and sample size before
      reducing evaluation rigor --- guided by the M0 measurements (Section~\ref{sec:m0}).
\item If long-tail diversification decreases human relevance materially, tune the exploration quota rather
      than forcing low-quality novelty.
\item If license/provenance evidence is contradictory, the system lowers confidence or flags the dataset
      (the \textsc{Conflicting} state, Section~\ref{sec:resolver}) rather than inventing a resolution.
\item If the full multi-modal pipeline becomes too large for Capstone~1, freeze one modality for the
      end-to-end experiment and retain the second modality as a validated extension.
\item \newtag If the reliability priors in Table~\ref{tab:reliability} produce poor calibration (high ECE)
      even after tuning, fall back to reporting raw evidence counts and source lists without a synthesized
      confidence label, rather than presenting an uncalibrated number as if it were meaningful.
\end{itemize}

% =========================================================================
\section{\newtag Worked Example: Evidence Resolution for a License Claim}
\label{app:worked-example}
% =========================================================================
Consider the claim $c = $ ``license of dataset $d$ is CC~BY~4.0,'' with three sources:

\begin{itemize}[leftmargin=*]
\item $s_1$: bare repository metadata field, asserts true, $r_1 = 0.60$.
\item $s_2$: official dataset card, asserts true, $r_2 = 0.70$.
\item $s_3$: an unverified community mirror README, asserts \emph{false} (claims ``All Rights Reserved''),
      $r_3 = 0.35$.
\end{itemize}

\textbf{Step 1 --- BPAs.}
$m_1(\top)=0.60,\ m_1(\Theta)=0.40$.
$m_2(\top)=0.70,\ m_2(\Theta)=0.30$.
$m_3(\bot)=0.35,\ m_3(\Theta)=0.65$.

\textbf{Step 2 --- Combine $m_1, m_2$.} No conflict yet since both support $\top$:
$m_{12}(\top) = m_1(\top)m_2(\top) + m_1(\top)m_2(\Theta) + m_1(\Theta)m_2(\top) = 0.60(0.70) + 0.60(0.30)
+ 0.40(0.70) = 0.42+0.18+0.28 = 0.88$.
$m_{12}(\Theta) = m_1(\Theta)m_2(\Theta) = 0.40(0.30) = 0.12$. $K_{12}=0$.

\textbf{Step 3 --- Combine with $m_3$.} Conflict mass:
$K = m_{12}(\top)\,m_3(\bot) = 0.88 \times 0.35 = 0.308$.
$m_{123}(\top) = \dfrac{m_{12}(\top)m_3(\Theta)}{1-K} = \dfrac{0.88 \times 0.65}{1-0.308} = \dfrac{0.572}{0.692}
\approx 0.827$.
$m_{123}(\bot) = \dfrac{m_{12}(\Theta)m_3(\bot)}{1-K} = \dfrac{0.12 \times 0.35}{0.692} = \dfrac{0.042}{0.692}
\approx 0.061$.
$m_{123}(\Theta) = \dfrac{m_{12}(\Theta)m_3(\Theta)}{1-K} = \dfrac{0.12 \times 0.65}{0.692} \approx 0.113$.

\textbf{Step 4 --- Interpret.} $\mathrm{Bel}(\top) \approx 0.827$, $\mathrm{Pl}(\top) \approx 0.940$,
$K_{\text{total}} \approx 0.308$. Since $K_{\text{total}} = 0.308 > \kappa_{\text{high}} = 0.30$, Algorithm~1
returns \textsc{Conflicting} \emph{despite} $\mathrm{Bel}(\top)$ being high --- correctly surfacing that
one credible-enough source actively disputes the claim, rather than averaging the dispute away. This is
the behavior a simple weighted mean cannot reproduce: a naive average of $\{0.60,0.70,(1-0.35)\!=\!0.65\}$
would report $\approx 0.65$ and hide that the sources are in outright disagreement rather than moderately
uncertain.

% =========================================================================
\section{\newtag Changelog: Version 1 to Version 2}
% =========================================================================
\begin{table}[h]
\centering
\small
\begin{tabular}{@{}p{1.1cm}p{5.6cm}p{5.6cm}@{}}
\toprule
\textbf{\#} & \textbf{Gap identified in review} & \textbf{Resolution in Version 2} \\
\midrule
1 & Evidence states were assigned heuristically, no combination rule. &
    Dempster--Shafer resolver, Algorithm~1, Section~\ref{sec:resolver}. \\
2 & Bias/diversity claims were asserted, not measured. &
    Exposure Disparity Index and calibration metrics, Section~\ref{sec:bias-calibration}. \\
3 & No comparison against the ``just ask an LLM'' alternative. &
    LLM zero-shot baseline added as ablation \#1, Section~\ref{sec:ablation}. \\
4 & No defined behavior when a dataset cannot be sampled at all. &
    Transfer-estimate fallback, Section~\ref{sec:coldstart}. \\
5 & Trust/explanation quality was a design principle, not evaluated. &
    Human trust study, Section~\ref{sec:human-eval}. \\
6 & M1-Air compute feasibility of the 30--100 candidate pool was assumed. &
    Milestone M0 validation spike, Section~\ref{sec:m0}. \\
7 & No task ownership across the 5--6 person team. &
    Subsystem ownership table, Section~\ref{sec:raci}. \\
8 & Automated sampling had no explicit ToS/compliance step. &
    Legal/ToS compliance checklist, Section~\ref{sec:legal}. \\
9 & CLI/demo interface was an optional extension (M6). &
    Moved into core MVP scope, Section~21.1. \\
\bottomrule
\end{tabular}
\end{table}

\textbf{Final project definition.} This system is a local, research-oriented dataset intelligence
pipeline that retrieves candidate datasets, resolves their evidence through a formal, testable fusion
rule, estimates task-specific usefulness (with an explicit fallback when direct measurement is
infeasible), measurably mitigates popularity bias, diversifies recommendations, calibrates its own
confidence, and validates its explanations against human judgment. Retrieval remains necessary, but
retrieval is no longer the endpoint --- and the system's claims about evidence, bias, and trust are now
each individually falsifiable.

\end{document}
